\documentclass[11pt,a4paper]{article}
\usepackage[T1]{fontenc}
\usepackage[utf8]{inputenc}
\usepackage[provide=*,lithuanian]{babel}
\usepackage{lmodern,amsmath,amssymb,enumitem}
\usepackage[margin=1.5cm]{geometry}
\setlength{\parindent}{0pt}
\setlength{\emergencystretch}{3em}
\setlist[enumerate]{label=\textbf{\arabic*.},leftmargin=0.7cm,itemsep=0.3cm,topsep=0.2cm,parsep=0pt}
\begin{document}
\begin{center}
{\Large\bfseries Vijetos teoremos taikymas}

\medskip
{\large Lygis: Baby}
\end{center}
\textbf{Užduotis.} Atlikite uždavinius, naudodamiesi Vijetos teorema.
\medskip
\begin{enumerate}
\item Duota lygtis \(x^2+5x+4=0\). Jos šaknys yra \(x_1,x_2\). Nerasdami šaknų, apskaičiuokite \mbox{\(x_1^2+x_2^2\)}.
\item Duotas daugianaris \(P(x)=x^2-7x+10\). Jis turi dvi realiąsias šaknis. Naudodamiesi Vijetos teorema, nustatykite šaknų ženklus, jų neskaičiuodami.
\item Viena daugianario \(P(x)=x^2+px-15\) šaknis yra \(3\). Raskite \(p\) ir kitą šaknį.
\item Sudarykite kvadratinį daugianarį \(P(x)\), kurio šaknys yra \(-2\) ir \(5\), o vyriausiasis koeficientas lygus \(1\).
\item Duotas daugianaris \(P(x)=x^2-5x+6\). Jo šaknys yra \(x_1,x_2\). Nerasdami šaknų, apskaičiuokite \mbox{\(\frac1{x_1}+\frac1{x_2}\)}.
\item Viena daugianario \(P(x)=x^2-7x+12\) šaknis yra \(3\). Raskite kitą šaknį, nespręsdami lygties \(P(x)=0\).
\item Duota lygtis \(x^2-6x+8=0\). Jos šaknys yra \(x_1,x_2\). Raskite \(x_1+x_2\) ir \(x_1x_2\).
\item Duotas daugianaris \(P(x)=x^2-9x+20\). Raskite jo šaknis, naudodamiesi Vijetos teorema.
\item Duotas daugianaris \(P(x)=x^2-7x+10\). Jo šaknys yra \(x_1,x_2\). Nerasdami šaknų, apskaičiuokite \mbox{\((x_1-1)(x_2-1)\)}.
\item Duota lygtis \(x^2-8x+12=0\). Jos šaknys yra \(x_1,x_2\). Nerasdami šaknų, apskaičiuokite \mbox{\((x_1-x_2)^2\)}.
\item Sudarykite kvadratinę lygtį, kurios šaknys yra \(3\) ir \(4\), o koeficientas prie \(x^2\) lygus \(1\).
\item Viena lygties \(x^2-3x-10=0\) šaknis yra \(-2\). Raskite kitą šaknį, nespręsdami lygties.
\item Duota lygtis \(x^2+3x-4=0\). Jos šaknys yra \(x_1,x_2\). Nerasdami šaknų, apskaičiuokite \mbox{\((x_1+2)(x_2+2)\)}.
\item Viena lygties \(x^2-7x+q=0\) šaknis yra \(2\). Raskite \(q\) ir kitą šaknį.
\item Duotas daugianaris \(P(x)=2x^2-10x+8\). Jo šaknys yra \(x_1,x_2\). Raskite \(x_1+x_2\) ir \(x_1x_2\).
\item Duota lygtis \(2x^2+3x-2=0\). Jos šaknys yra \(x_1,x_2\). Nerasdami šaknų, apskaičiuokite \mbox{\(\frac1{x_1}+\frac1{x_2}\)}.
\item Duotas daugianaris \(P(x)=x^2+6x+7\). Jo šaknys yra \(x_1,x_2\). Nerasdami šaknų, apskaičiuokite \mbox{\((x_1-x_2)^2\)}.
\item Duota lygtis \(x^2-4x-12=0\). Išspręskite ją, naudodamiesi Vijetos teorema.
\item Duotas daugianaris \(P(x)=2x^2-6x+1\). Jo šaknys yra \(x_1,x_2\). Nerasdami šaknų, apskaičiuokite \mbox{\(x_1^2+x_2^2\)}.
\item Duota lygtis \(2x^2+14x+12=0\). Ji turi dvi realiąsias šaknis. Naudodamiesi Vijetos teorema, nustatykite šaknų ženklus, jų neskaičiuodami.
\item Viena daugianario \(P(x)=2x^2+4x-6\) šaknis yra \(1\). Raskite kitą šaknį, nespręsdami lygties \(P(x)=0\).
\item Duota lygtis \(-x^2+3x+10=0\). Jos šaknys yra \(x_1,x_2\). Raskite \(x_1+x_2\) ir \(x_1x_2\).
\item Duotas daugianaris \(P(x)=-x^2+6x-8\). Jo šaknys yra \(x_1,x_2\). Nerasdami šaknų, apskaičiuokite \mbox{\(\frac1{x_1}+\frac1{x_2}\)}.
\item Sudarykite kvadratinį daugianarį \(P(x)\), kurio šaknys yra \(-1\) ir \(-6\), o vyriausiasis koeficientas lygus \(1\).
\item Duotas daugianaris \(P(x)=2x^2-6x+1\). Jo šaknys yra \(x_1,x_2\). Nerasdami šaknų, apskaičiuokite \mbox{\((x_1-3)(x_2-3)\)}.
\item Duotas daugianaris \(P(x)=-x^2-7x-12\). Raskite jo šaknis, naudodamiesi Vijetos teorema.
\item Duotas daugianaris \(P(x)=x^2-x-12\). Jis turi dvi realiąsias šaknis. Naudodamiesi Vijetos teorema, nustatykite šaknų ženklus, jų neskaičiuodami.
\item Duota lygtis \(2x^2-4x-1=0\). Jos šaknys yra \(x_1,x_2\). Nerasdami šaknų, apskaičiuokite \mbox{\((x_1-x_2)^2\)}.
\item Viena daugianario \(P(x)=2x^2+px-6\) šaknis yra \(1\). Raskite \(p\) ir kitą šaknį.
\item Duota lygtis \(-x^2+4x+2=0\). Jos šaknys yra \(x_1,x_2\). Nerasdami šaknų, apskaičiuokite \mbox{\(x_1^2+x_2^2\)}.
\item Duota lygtis \(-x^2+2x+3=0\). Jos šaknys yra \(x_1,x_2\). Nerasdami šaknų, apskaičiuokite \mbox{\((2x_1+1)(2x_2+1)\)}.
\item Viena lygties \(-x^2+5x+6=0\) šaknis yra \(-1\). Raskite kitą šaknį, nespręsdami lygties.
\item Duota lygtis \(x^2+2x-3=0\). Jos šaknys yra \(x_1,x_2\). Nerasdami šaknų, apskaičiuokite \mbox{\(\frac1{x_1}+\frac1{x_2}\)}.
\item Duotas daugianaris \(P(x)=-x^2-3x+1\). Jo šaknys yra \(x_1,x_2\). Nerasdami šaknų, apskaičiuokite \mbox{\((x_1-x_2)^2\)}.
\item Sudarykite kvadratinę lygtį, kurios šaknys yra \(-1\) ir \(3\), o koeficientas prie \(x^2\) lygus \(2\).
\item Duotas daugianaris \(P(x)=x^2-3x-2\). Jo šaknys yra \(x_1,x_2\). Nerasdami šaknų, apskaičiuokite \mbox{\(x_1^2+x_2^2\)}.
\item Duotas daugianaris \(P(x)=3x^2+3x-6\). Jo šaknys yra \(x_1,x_2\). Raskite \(x_1+x_2\) ir \(x_1x_2\).
\item Duota lygtis \(x^2+5x=0\). Ji turi dvi realiąsias šaknis. Naudodamiesi Vijetos teorema, nustatykite šaknų ženklus, jų neskaičiuodami.
\item Duota lygtis \(x^2+3x-10=0\). Išspręskite ją, naudodamiesi Vijetos teorema.
\item Viena lygties \(-x^2+5x+q=0\) šaknis yra \(-1\). Raskite \(q\) ir kitą šaknį.
\item Duota lygtis \(-2x^2-10x-8=0\). Jos šaknys yra \(x_1,x_2\). Raskite \(x_1+x_2\) ir \(x_1x_2\).
\item Duotas daugianaris \(P(x)=2x^2-7x+3\). Raskite jo šaknis, naudodamiesi Vijetos teorema.
\item Viena daugianario \(P(x)=3x^2-5x-2\) šaknis yra \(2\). Raskite kitą šaknį, nespręsdami lygties \(P(x)=0\).
\item Viena daugianario \(P(x)=x^2+px+12\) šaknis yra \(-3\). Raskite \(p\) ir kitą šaknį.
\item Duotas daugianaris \(P(x)=x^2-4x+2\). Jo šaknys yra \(x_1,x_2\). Nerasdami šaknų, apskaičiuokite \mbox{\((x_1+1)(x_2+1)\)}.
\item Duotas daugianaris \(P(x)=3x^2-7x+2\). Jo šaknys yra \(x_1,x_2\). Nerasdami šaknų, apskaičiuokite \mbox{\(\frac1{x_1}+\frac1{x_2}\)}.
\item Duota lygtis \(3x^2+6x-1=0\). Jos šaknys yra \(x_1,x_2\). Nerasdami šaknų, apskaičiuokite \mbox{\(x_1^2+x_2^2\)}.
\item Sudarykite kvadratinį daugianarį \(P(x)\), kurio šaknys yra \(-2\) ir \(4\), o vyriausiasis koeficientas lygus \(-1\).
\item Duota lygtis \(x^2+x-5=0\). Jos šaknys yra \(x_1,x_2\). Nerasdami šaknų, apskaičiuokite \mbox{\((x_1-x_2)^2\)}.
\item Duotas daugianaris \(P(x)=-x^2-7x-10\). Jis turi dvi realiąsias šaknis. Naudodamiesi Vijetos teorema, nustatykite šaknų ženklus, jų neskaičiuodami.
\end{enumerate}
\clearpage
\section*{Atsakymai}
\begingroup\small
\setlength{\baselineskip}{11pt}
\setlist[enumerate]{itemsep=0cm}
\begin{enumerate}
\item \(17\).
\item Abi šaknys teigiamos.
\item \(p=2,\quad x_2=-5\).
\item \(P(x)=x^2-3x-10\).
\item \(\frac{5}{6}\).
\item \(x_2=4\).
\item \(x_1+x_2=6,\quad x_1x_2=8\).
\item \(x_1=4,\quad x_2=5\).
\item \(4\).
\item \(16\).
\item \(x^2-7x+12=0\).
\item \(x_2=5\).
\item \(-6\).
\item \(q=10,\quad x_2=5\).
\item \(x_1+x_2=5,\quad x_1x_2=4\).
\item \(\frac{3}{2}\).
\item \(8\).
\item \(x_1=-2,\quad x_2=6\).
\item \(8\).
\item Abi šaknys neigiamos.
\item \(x_2=-3\).
\item \(x_1+x_2=3,\quad x_1x_2=-10\).
\item \(\frac{3}{4}\).
\item \(P(x)=x^2+7x+6\).
\item \(\frac{1}{2}\).
\item \(x_1=-4,\quad x_2=-3\).
\item Viena šaknis teigiama, kita neigiama.
\item \(6\).
\item \(p=4,\quad x_2=-3\).
\item \(20\).
\item \(-7\).
\item \(x_2=6\).
\item \(\frac{2}{3}\).
\item \(13\).
\item \(2x^2-4x-6=0\).
\item \(13\).
\item \(x_1+x_2=-1,\quad x_1x_2=-2\).
\item Viena šaknis lygi \(0\), kita neigiama.
\item \(x_1=-5,\quad x_2=2\).
\item \(q=6,\quad x_2=6\).
\item \(x_1+x_2=-5,\quad x_1x_2=4\).
\item \(x_1=\frac{1}{2},\quad x_2=3\).
\item \(x_2=-\frac{1}{3}\).
\item \(p=7,\quad x_2=-4\).
\item \(7\).
\item \(\frac{7}{2}\).
\item \(\frac{14}{3}\).
\item \(P(x)=-x^2+2x+8\).
\item \(21\).
\item Abi šaknys neigiamos.
\end{enumerate}
\endgroup
\end{document}
